\label{app:emissionrule}

Every method compared here emits a number of candidates fixed by convention rather than by the
substrate: a budget the field reports at, or whatever a threshold admits. The references are not
shaped that way. Enzymes compete for one substrate, so its annotated products are alternatives drawn
from a shared pool, and how many there are varies: the median substrate on the clean test split has
one and the mean $2.22$, against a deployed policy that emits $8.4$ for every substrate alike.

This appendix reports what that costs, with the ranking held fixed in every arm so that any
difference is attributable to the size of the emitted set and to nothing else. Three quantities
bound the question. The deployed budget is what the system ships. The best global constant is the
best any budget-based policy could have done, chosen by looking at this split. Truncating each
substrate at its own annotated count is the oracle: it is what perfect knowledge of the set size
would be worth at this ranking quality, and it is not achievable.

\paragraph{The comparison that decides the claim.} A relative rule emits while a candidate scores
within a factor $\alpha$ of the leader, which is what a model normalised over its siblings does at
inference. Against the deployed budget it gains $+0.074$ $[+0.059,+0.089]$ in macro F1, but that is
the weak arm and the paper does not rest on it: the deployed budget is a policy this appendix is
arguing against, and beating it shows only that fifteen is too many. The arm that decides whether
the \emph{pool} is the right object is the best global constant, because a constant is precisely the
policy that ignores the pool. Both arms are chosen with the same hindsight on the same split, which
is stated rather than hidden, since neither is a forecast.

\begin{center}\small
\begin{tabular}{lcccl}
\toprule
policy & macro F1 & over the best constant & 95\% CI & separated \\
\midrule
deployed budget, $k=15$        & $0.136$ & $-0.059$ & --- & --- \\
best global constant, $k=1$    & $0.195$ & --- & --- & --- \\
\midrule
relative rule, $\alpha=0.8$    & $0.201$ & $+0.0061$ & $[+0.0012,+0.0113]$ & yes \\
relative rule, $\alpha=0.75$   & $0.202$ & $+0.0074$ & $[+0.0019,+0.0129]$ & yes \\
relative rule, $\alpha=0.6$    & $0.206$ & $+0.0114$ & $[+0.0042,+0.0188]$ & yes \\
\textbf{relative rule, $\alpha=0.5$} & $\mathbf{0.209}$ & $\mathbf{+0.0147}$ & $[+0.0059,+0.0237]$ & yes \\
relative rule, $\alpha=0.4$    & $0.206$ & $+0.0108$ & $[+0.0015,+0.0204]$ & yes \\
relative rule, $\alpha=0.3$    & $0.206$ & $+0.0113$ & $[+0.0005,+0.0221]$ & yes \\
\midrule
oracle: each substrate's own count & $0.214$ & $+0.0189$ & $[+0.0126,+0.0256]$ & yes \\
forecast of that count             & $0.181$ & $-0.0140$ & $[-0.0262,-0.0015]$ & yes, worse \\
\bottomrule
\end{tabular}
\end{center}

\paragraph{The threshold is not fitted.} A rule that works at one value of a free parameter is a fit
to the split it was tuned on. Sweeping $\alpha$ over ten values from $0.95$ to $0.1$ gives a
plateau, not a spike: six consecutive values from $0.8$ down to $0.3$ beat the best constant with an
interval excluding zero, the gain rises to $\alpha=0.5$ and falls away symmetrically, and the two
extremes fail in the two ways they should. At $\alpha=0.95$ the rule is a constant of one and gains
$+0.0013$ $[-0.0007,+0.0036]$; at $\alpha=0.1$ it admits nearly everything and loses. What the
sweep shows is that the pool carries information about how many to emit across a wide band of
thresholds, which is the claim. But $0.5$ is an argmax over the sweep, which is a
tuned value whatever the plateau around it looks like, so it is cross-fitted rather than defended.
Choosing both $\alpha$ and the competing constant on four fifths of the substrates and comparing
only on the remaining fifth, over $100$ random splits, reproduces the point estimate at $+0.0151$
and selects $\alpha=0.5$ as the training-part argmax in all $100$. The spread across held-out fifths
is wide, $[-0.0015,+0.0334]$, as a fifth of a thousand substrates must give; what this establishes
is that the threshold is not a fit to this split, and it is not a second interval on the gain.

A companion sweep over \emph{absolute} thresholds is released beside it
(\texttt{results/stopping\_rule.json}) and is a negative result: cutting the pool at a fixed score
gains $+0.0035$ $[-0.0047,+0.0120]$ and calibrating it $+0.0039$ $[-0.0042,+0.0122]$, neither
separated from zero. Only the two pool-relative rules separate --- the one reported here and an
expected-F1 rule at $+0.0126$ $[+0.0041,+0.0214]$. That is the comparison the claim needs: the
information is in where a candidate sits relative to the pool, not in its score.

\paragraph{Four fifths of the oracle, without the oracle.} The relative rule reaches $+0.0147$ of
an available $+0.0189$. That is the useful form of the result: the headroom over the best constant
is small in absolute terms, and almost all of it is reachable without knowing anything the evaluator
knows. Forecasting the count directly from the substrate, in contrast, is worse than the constant it
replaces by an interval that excludes zero, so the count is not predictable from the substrate in
the way the pool ordering is (Appendix~\ref{app:negative}).

\paragraph{Where it does not hold.} The advantage is not uniform across the split's two populations
(\S\ref{sec:robust}). On the $845$ parent drugs the rule beats the best constant by $+0.0210$
$[+0.0101,+0.0316]$; on the $325$ substrates that are themselves metabolites of another substrate
here it does not, at $-0.0016$ $[-0.0160,+0.0126]$, so on the stratum where the model is weaker its
scores say nothing about how many to emit. Under strict \textsc{inchikey} matching rather than the
tautomer-aware default the margin is $+0.0151$ $[+0.0063,+0.0241]$, so the result is not an artifact
of the most permissive criterion. It remains a property of one method's frozen scores on one corpus:
it says that this pool is informative about this set size, not that relative emission is generally
preferable.
